\subsection{问题二模型的建立与求解}

\subsubsection{排程对象与完成口径}

问题二只加载住院事件，不预占任何门诊、体检容量。记$p_i$为医嘱下达时刻，$r_i$为$p_i$向上取整到5分钟后在工作日历中的最早可服务时刻；48小时截止时刻为
\begin{equation}
d_i=p_i+48\unit{h}.
\end{equation}
若所有$j\in\J_i$均在$d_i$前结束，则$y_i=1$，否则$y_i=0$。住院48小时完整完成率为
\begin{equation}
R_I=\frac{\sum_{i\in\I_I^{\mathrm{eval}}}y_i}{|\I_I^{\mathrm{eval}}|}.
\end{equation}
留出期评价集合包含70\,771个事件。258个3月30--31日预热事件可占用4月初资源，但既不进入分子也不进入分母；3月31日事件的截止期可延伸到4月2日，故日历尾端保留2日而不按月重置。

\subsubsection{时间索引整数规划}

令$x_{ijrt}=1$表示事件$i$的项目$j$在设备$r$、时隙$t$开始，$a_{ijrtu}=1$表示该选择会占用时隙$u$。时长以5分钟向上取整：
\begin{equation}
\ell_{ij}=\left\lceil\frac{\tau_{ij}}{5}\right\rceil.
\end{equation}
事件采用“全部安排或全部不安排”的完整性约束：
\begin{equation}
\sum_{r\in\R_{ij}^{\mathrm{final}}}\sum_{t\in\T_{ijr}}x_{ijrt}=y_i,
\qquad \forall i,\ j\in\J_i.
\end{equation}
由此，事件被选择时每个项目恰好安排一次，未被选择时不占用任何资源。

设备互斥约束为
\begin{equation}
\sum_i\sum_j\sum_{t:a_{ijrtu}=1}x_{ijrt}\le1,
\qquad \forall r,u.
\end{equation}
同一患者不能同时检查，且同一患者在不同时刻开出的多个申请事件仍共享同一患者标识$\pi(i)$：
\begin{equation}
\sum_{i:\pi(i)=v}\sum_j\sum_r\sum_{t:a_{ijrtu}=1}x_{ijrt}\le1,
\qquad \forall v,u.
\end{equation}
医生是所有设备共享的并发资源。设时隙$u$对应星期$w(u)$，训练期容量为$C_{w(u),u}$，则
\begin{equation}
\sum_i\sum_j\sum_r\sum_{t:a_{ijrtu}=1}x_{ijrt}\le C_{w(u),u},\qquad\forall u.
\end{equation}
这比“每个房间各有若干医生”的假设更符合题面，因为33名医生应在22台设备之间共享。

释放与截止通过可选开始集合$\T_{ijr}$保证：
\begin{align}
t&\ge r_i,\\
t+5\ell_{ij}&\le d_i.
\end{align}
任一项目不得跨越12:00--13:00午休或日界。对空腹项目，进一步要求
\begin{equation}
t+5\ell_{ij}\le 10{:}00.
\end{equation}
对憋尿项目，项目级最早开始为
\begin{equation}
t\ge \left\lceil p_i+60\unit{min}\right\rceil_{5\mathrm{min}}.
\end{equation}
床旁项目的设备集合已由$\R_{ij}\cap\{7\}$给出。

为使项目相对顺序成为决策，记$X_{ij}=\sum_{r,t}x_{ijrt}$。对任意两个不同任务$(i,j)$与$(k,\ell)$，若$\pi(i)=\pi(k)$，令$z_{ij,k\ell}=1$表示前者先于后者。仅当两个任务均被选中时，必须且只能选择一个方向：
\begin{align}
z_{ij,k\ell}+z_{k\ell,ij}&\le X_{ij},\\
z_{ij,k\ell}+z_{k\ell,ij}&\le X_{k\ell},\\
z_{ij,k\ell}+z_{k\ell,ij}&\ge X_{ij}+X_{k\ell}-1.
\end{align}
设备切换的转运间隔定义为
\begin{equation}
g(q_{ij},q_{k\ell})=
\begin{cases}
0,&q_{ij}=q_{ik},\\
10\unit{min},&q_{ij}\ne q_{ik}.
\end{cases}
\end{equation}
利用足够大的常数$M$关闭未选方向，成对非重叠与转运约束为
\begin{align}
s_{k\ell}&\ge s_{ij}+5\ell_{ij}+g(q_{ij},q_{k\ell})-M(1-z_{ij,k\ell}),\\
s_{ij}&\ge s_{k\ell}+5\ell_{k\ell}+g(q_{k\ell},q_{ij})-M(1-z_{k\ell,ij}).
\end{align}
时间索引变量决定开始时刻，顺序变量表达同一患者在事件内及跨事件的成对先后关系；构造算法则在安排每个任务时同步检查该患者已占用的全部区间。问题二的主目标为
\begin{equation}
\max\ \sum_{i\in\I_I^{\mathrm{eval}}}y_i.
\end{equation}
完成数相同时，再按等待时间和跨室等运营指标选择推荐策略，不用次级指标交换一个完整完成事件。对完整排入的事件，等待时间从最早可服务时刻计至首项检查开始：
\begin{equation}
w_i=\min_{j\in\J_i}s_{ij}-r_i,\qquad y_i=1.
\end{equation}
人员在岗容量利用率按全部已排任务分钟占可用医生分钟的比例计算：
\begin{equation}
U_{\mathrm{staff}}
=\frac{\sum_{i,j,r,t}5\ell_{ij}x_{ijrt}}
{5\sum_{u\in\T}C_{w(u),u}}.
\end{equation}
多项目跨室率以评价期全部多项目住院事件为分母；完整排入且至少发生一次设备切换的事件进入分子，未完整排入者不产生实际跨室。因此该指标用于比较患者移动负担时，需要与完整完成率同时考察。

\subsubsection{大规模构造算法}

全年时间索引模型过大，本文保留上述可行域，改用确定性患者级构造算法求解。诊断资源的动态优先级和住院患者柔性排程已有类似建模思想\cite{patrick2008,patrick2007}，本文算法内核由“候选序列—共同设备—最早可行插入”组成。

\textbf{步骤1：候选项目序列。} 对每个事件生成四类确定性序列：空腹优先后依次按兼容设备数少、时长长排序；空腹优先后依次按时长长、兼容设备数少排序；兼容设备数少后按空腹优先排序；以及原始项目序。重复序列只保留一次，所有序列均通过患者互斥约束执行。

\textbf{步骤2：共同设备优先。} 先计算
\begin{equation}
\R_i^{\cap}=\bigcap_{j\in\J_i}\R_{ij}^{\mathrm{final}}.
\end{equation}
若$\R_i^{\cap}\ne\varnothing$，先按相应设备规则排序，取前三个共同设备，并用“空腹优先、设备选择少者优先、长项目优先”的首选项目序列尝试同室完成；这些尝试均失败后，再枚举其余候选序列并允许跨室。所有策略均先选取最早可行时刻；当多个设备同样早时，分别比较设备编号、当前累计负荷或训练期稀缺权重
\begin{equation}
\omega_r=\sum_j\frac{D_j\tau_j}{|\R_j|}\indicator(r\in\R_j),
\end{equation}
由此形成最低编号、负荷均衡和稀缺能力保留三种设备破同序方式，最后一种用于减少普通任务对高需求专项设备的占用。

\textbf{步骤3：连续日历搜索。} 对候选序列逐任务向前搜索。某一插入位置可行，当且仅当设备区间空闲、每个5分钟槽的医生负载未达上限、项目在单一工作时段内、准备条件满足且最终不越过$d_i$。若有多个位置，先选择开始最早者；同样早时优先延续上一设备，再按相应候选策略的设备编号、累计负荷或稀缺权重破同序。

\textbf{步骤4：整体可行性确认。} 只有所有项目都找到位置，才确认该事件占用的设备与医生时段；中途失败则取消本次尝试的全部暂存位置，不保留部分任务。该规则与患者完整完成口径一致。

\textbf{步骤5：患者优先级。} 令$\xi_i=1$表示事件$i$被抽为题设复杂病例，主情景总工时为
\begin{equation}
W_i^{\star}=(1-\xi_i)W_i^{(0)}+\xi_iW_i^{(1)}.
\end{equation}
比较两种住院队列：FCFS按$r_i$和事件编号排序；最小松弛度保护按
\begin{equation}
\mathrm{slack}_i=(d_i-p_i)-W_i^{\star}
\end{equation}
升序，并依次用$r_i$、候选设备数、截止时刻和事件编号破同序。$d_i-p_i$表示从医嘱下达到截止的可用时长。两种患者顺序分别与最低编号、负荷均衡和稀缺能力保留三种设备规则组合，共形成六个候选策略；六者使用相同的项目序列、准备条件与可行性内核。对策略集合$\mathcal S_2$，固定选择规则为
\begin{equation}
\label{eq:p2-policy-selection}
\begin{aligned}
S_2^{\star}=\arg\min_{S\in\mathcal S_2}\bigl(&-N_I(S),P_{90}(S),P_{50}(S),\\
&r_{\mathrm{switch}}(S),M(S),\mathrm{id}(S)\bigr).
\end{aligned}
\end{equation}
其中$N_I$为住院完整完成数，$M$为排入总分钟，最后一项仅用于稳定破同序。负号保证任何次级指标都不以减少一个完整事件为代价。

\subsubsection{问题二求解结果}

表\ref{tab:p2-policy}比较两种患者队列与三种设备破同序方式组成的六个策略，图\ref{fig:p2-policy}进一步对照其完成率与条件等待。FCFS--负荷均衡完整完成60\,368个留出期住院事件，完成率85.30\%，在六个策略中最高；其条件等待P50和P90分别为3.58小时和22.08小时。按照式（\ref{eq:p2-policy-selection}）的词典序规则，推荐FCFS--负荷均衡策略。表中多项目需求跨室率采用前述固定分母，因而在六个策略间的排序等价于比较发生跨室的事件数。对比结果还表明，设备选择规则与患者队列共同影响可完成集合，不能只根据FCFS或松弛度标签判断策略优劣。

\begin{table}[H]
\centering
\caption{问题二纯住院策略比较}
\label{tab:p2-policy}
\begin{tabular}{lrrrrr}
\toprule
策略 & 48h完成数 & 完成率 & P50/h & P90/h & 多项目需求跨室率 \\
\midrule
FCFS--最低编号 & 59,586 & 84.20\% & 25.50 & 47.50 & 57.85\% \\
最小松弛度--最低编号 & 58,964 & 83.32\% & 20.25 & 46.00 & 60.76\% \\
FCFS--负荷均衡 & 60,368 & 85.30\% & 3.58 & 22.08 & 58.16\% \\
最小松弛度--负荷均衡 & 58,419 & 82.55\% & 3.33 & 21.67 & 57.29\% \\
FCFS--稀缺能力保护 & 60,184 & 85.04\% & 6.08 & 29.42 & 57.83\% \\
最小松弛度--稀缺能力保护 & 58,875 & 83.19\% & 6.42 & 29.67 & 58.08\% \\
\bottomrule
\end{tabular}

\end{table}

\begin{figure}[H]
\centering
\includegraphics[width=.92\textwidth]{fig04_p2_policy.pdf}
\caption{问题二住院完成率与条件等待比较}
\label{fig:p2-policy}
\end{figure}

在评价分母中，10\,403个事件未能完整排入。表\ref{tab:p2-failure}按最终失败状态作互斥分解：10\,258个事件属于项目设备能力不满足，其中222个事件含有表1无法确认设备的项目，10\,036个事件的床旁任务在“项目兼容设备”与7号床旁机取交集后为空，其他A--F能力失败类别均为0；另有145个事件虽各项目存在候选设备，但在48小时、准备窗口和已有负荷共同作用下没有形成完整计划。前一结果说明床旁地点能力与专项检查能力的交叉覆盖是主要瓶颈，不能用空闲普通诊室替代；医院若核定7号机具有更多专项能力，应先更新项目资质表再重新排程。

\begin{table}[H]
\centering
\caption{问题二评价期未完成原因分解}
\label{tab:p2-failure}
\begin{tabular}{lrr}
\toprule
评价期未完成原因 & 事件数 & 占未完成事件 \\
\midrule
无满足项目能力的设备组合 & 10,245 & 98.40\% \\
48小时内未形成完整可行计划 & 167 & 1.60\% \\
\midrule
合计 & 10,412 & 100.00\% \\
\bottomrule
\end{tabular}

\end{table}

完整完成者的P90等待为22.08小时、P95为23.67小时。多项目患者中58.16\%发生跨室，这一比例较高，原因是不同项目的专项设备集合交集不足；相邻项目更换实际诊室时均计入10分钟转运。

纯住院方案共排入1\,344\,765分钟，人员在岗容量利用率为54.53\%。这一总体利用率不能说明各机器均有富余：表\ref{tab:p2-room-load}中排入分钟最高的5个诊室，其有任务日平均利用率约为54.96\%--54.97\%，负荷均衡规则已使普通兼容任务较均匀地分布；与此同时，能力集合为空的任务即使遇到普通机器空闲也不能转入。床旁、血管、心脏和产科任务仍受各自候选集合限制。因此问题二的主要改进方向是核定真实项目资质并均衡相容任务，而不是只提高全科总开机率。

\begin{table}[H]
\centering
\caption{问题二排入分钟最高的诊室}
\label{tab:p2-room-load}
\begin{tabular}{lrr}
\toprule
诊室 & 排入分钟 & 有任务日平均利用率 \\
\midrule
门6 & 96,840 & 54.97\% \\
门11 & 96,840 & 54.97\% \\
门7 & 96,830 & 54.97\% \\
病4 & 96,825 & 54.96\% \\
门5 & 96,825 & 54.96\% \\
门10 & 55,700 & 31.62\% \\
\bottomrule
\end{tabular}

\end{table}

问题二输出患者、项目、开始时刻、结束时刻和设备组成的推荐表，部分结果见表\ref{tab:recommendation}。每个项目独立保留，事件内序号给出实际执行顺序；相邻项目跨设备时，下一项目开始前已经计入转运。

\begin{table}[H]
\centering
\caption{患者级推荐排程样例}
\label{tab:recommendation}
\begin{tabularx}{\textwidth}{r l r X l l}
\toprule
序位 & 事件 & 项目序位 & 项目族 & 时段 & 设备 \\
\midrule
1 & IP2024-000259 & 1 & 腹部/胃肠 & 04-01 08:30--08:40 & 门16(65) \\
2 & IP2024-000261 & 1 & 浅表器官 & 04-01 08:30--08:40 & 妇门2(75) \\
3 & IP2024-000265 & 1 & 浅表器官 & 04-01 08:35--08:45 & 门9(47) \\
4 & IP2024-000268 & 1 & 血管 & 04-01 08:50--09:00 & 病2(71) \\
4 & IP2024-000268 & 2 & 血管 & 04-01 09:40--09:50 & 病2(71) \\
5 & IP2024-000267 & 1 & 浅表器官 & 04-01 09:00--09:10 & 病3(72) \\
6 & IP2024-000271 & 1 & 血管 & 04-01 09:00--09:10 & 门3(42) \\
6 & IP2024-000271 & 2 & 介入/定位 & 04-01 09:45--09:55 & 门3(42) \\
\bottomrule
\end{tabularx}

\end{table}

综合48小时完成数、等待时间和跨室比例，最终选取FCFS--负荷均衡作为住院策略。患者内部先处理候选设备较少和准备窗口较紧的项目，优先使用共同设备，无法同室完成时再按最早可行时刻跨室安排。推荐表由此同时给出患者顺序、项目次序、检查时刻和诊室；全年构造解的求解质量由第六章的精确子问题对照界定。
